\begin{thema}{A}
\begin{erwthma}
    \item Να διατυπώσετε και να δώσετε τη γεωμετρική ερμηνεία του θεωρήματος \eng{Rolle}.
    \item Τι ονομάζουμε αρχική συνάρτηση ή παράγουσα μιας συνάρτησης $f$ σε ένα διάστημα $\Delta$?
    \item Πώς ορίζονται οι πράξεις του γινομένου $f\cdot g$ και του πηλίκου $\dfrac{f}{g}$ δύο συναρτήσεων? Ποιο το πεδίο ορισμού τους?
    \item Να χαρακτηρίσετε τις προτάσεις που ακολουθούν, γράφοντας στο τετράδιό σας, δίπλα στο γράμμα που αντιστοιχεί σε κάθε πρόταση, τη λέξη \textbf{Σωστό}, αν η πρόταση είναι σωστή, ή \textbf{Λάθος}, αν η πρόταση είναι λανθασμένη.
    \begin{alist}
        \item Αν $\lim\limits_{x \to x_0} f(x) = 0$ και η $g$ είναι φραγμένη κοντά στο $x_0$, τότε $\lim\limits_{x \to x_0} [f(x)g(x)] = 0$.
        \item Αν δύο συναρτήσεις $f, g$ έχουν την ίδια παράγωγο στο $\Delta$, τότε οι εφαπτομένες των $C_f, C_g$ στα σημεία με μια κοινή τετμημένη, είναι είτε παράλληλες είτε ταυτίζονται.
        \item Το $\dint_a^\beta f(x) \d x$ είναι ένας πραγματικός αριθμός, που εξαρτάται από τα όρια $a, \beta$ και τη συνάρτηση.
        \item Αν $\lim\limits_{x\to x_0} f(x) > 0$, τότε υποχρεωτικά $f(x) > 0$ για κάθε $x$ κοντά στο $x_0$.
        \item Αν $\lim\limits_{x\to x_0^+} f(x) \ne \lim\limits_{x\to x_0^-} f(x)$, τότε η $f$ δεν είναι συνεχής στο $x_0$.
    \end{alist}
    \item Αν $\dint_a^\beta f(x) \d x = 0$ και η $f$ είναι συνεχής με $f(x) \ge 0$, τότε αναγκαστικά:
    \begin{tasks}[label=\let\textdexiakeraia\relax\alph*.](2)
        \task $f(x) = 0$ για κάθε $x \in [a, \beta]$.
        \task H $f$ είναι σταθερή αλλά όχι μηδενική.
        \task $a = 0$ ή $\beta = 0$.
        \task H $f$ είναι περιττή.
    \end{tasks}
\end{erwthma}
\end{thema}
